\scanpage{55}
\mspage{27}
La première assertion s'obtient en ``bouchant'' le trou $i$, alors la section envisagée devient la section de $U\cup\{i\}=U'$ associée au pt $i$, et elle est donc distincte de celle associée aux \deleted{pts} autres pts de $U'$, i.e. aux pts de $U$ --- a fortiori avant de passer au quotient par le ss-groupe inv. engendré par $L_i$. On \unsure{voudra} de $\unsure{=|\overline m}$ façon se convaincre que les classes de conj. des scindages associés à un $L_{i_1}$ et un $L_{i_2}$, $i_1\ne i_2$, sont distinctes --- ça marche dans tous les cas sauf le cas du type $(0,3)$, --- car en bouchant deux trous, on tombe sur le type $(0,1)$, où on ne dispose pas du résultat d'injectivité. Mais on peut s'en tirer autrement, à condition d'admettre que pour un scindage de (2), faisant opérer $\Gamma$ sur $\pi$, on a
\[
(3)\qquad\pi^{\Gamma^\circ}=L_i
\]
(donc $\pi^\Gamma=(1)$, d'ailleurs) --- résultat qui me paraît hautement plausible. Cela
\margin{\figure{\begin{tikzpicture}\draw (0,0)--(0,-1.9);\draw (0.12,-0.5)--(0.12,-1.8);\end{tikzpicture}}}
\margin{À énoncer parmi les conjectures-clefs\ldots}
\margin{54}
